\documentclass[11pt,a4paper]{article}

% -------------------------------------------------
% Packages
% -------------------------------------------------
\usepackage[margin=2.2cm]{geometry}
\usepackage{kotex}
\usepackage{amsmath,amssymb}
\usepackage{booktabs}
\usepackage{array}
\usepackage{xcolor}
\usepackage{listings}
\usepackage[most]{tcolorbox}
\usepackage{enumitem}
\usepackage{hyperref}
\usepackage{fancyhdr}
\usepackage{titlesec}

% -------------------------------------------------
% Colors
% -------------------------------------------------
\definecolor{codebg}{RGB}{247,247,247}
\definecolor{codeframe}{RGB}{210,210,210}
\definecolor{keywordcolor}{RGB}{0,70,160}
\definecolor{commentcolor}{RGB}{0,120,80}
\definecolor{stringcolor}{RGB}{170,50,50}
\definecolor{titleblue}{RGB}{35,75,130}

% -------------------------------------------------
% Hyperlinks
% -------------------------------------------------
\hypersetup{
    colorlinks=true,
    linkcolor=titleblue,
    urlcolor=titleblue
}

% -------------------------------------------------
% Header / Footer
% -------------------------------------------------
\pagestyle{fancy}
\fancyhf{}
\lhead{C++ Programming}
\rhead{Day 1}
\cfoot{\thepage}

% -------------------------------------------------
% Section style
% -------------------------------------------------
\titleformat{\section}
  {\Large\bfseries\color{titleblue}}
  {\thesection.}{0.6em}{}

\titleformat{\subsection}
  {\large\bfseries}
  {\thesubsection.}{0.5em}{}

% -------------------------------------------------
% C++ code style
% -------------------------------------------------
\lstdefinestyle{cppstyle}{
    language=C++,
    backgroundcolor=\color{codebg},
    frame=single,
    rulecolor=\color{codeframe},
    basicstyle=\ttfamily\small,
    keywordstyle=\color{keywordcolor}\bfseries,
    commentstyle=\color{commentcolor},
    stringstyle=\color{stringcolor},
    numbers=left,
    numberstyle=\tiny\color{gray},
    numbersep=8pt,
    tabsize=4,
    showstringspaces=false,
    breaklines=true,
    columns=fullflexible
}

\lstset{style=cppstyle}

% -------------------------------------------------
% Boxes
% -------------------------------------------------
\newtcolorbox{learningbox}{
    colback=blue!4,
    colframe=titleblue,
    title=\textbf{학습 목표},
    fonttitle=\bfseries
}

\newtcolorbox{importantbox}{
    colback=yellow!8,
    colframe=orange!70!black,
    title=\textbf{핵심 포인트},
    fonttitle=\bfseries
}

\newtcolorbox{checkpointbox}{
    colback=green!5,
    colframe=green!45!black,
    title=\textbf{체크 질문},
    fonttitle=\bfseries
}

\newtcolorbox{practicebox}{
    colback=red!3,
    colframe=red!55!black,
    title=\textbf{실습},
    fonttitle=\bfseries
}

% -------------------------------------------------
% Document
% -------------------------------------------------
\begin{document}

\begin{titlepage}
    \centering
    \vspace*{3cm}

    {\Huge\bfseries C++ Programming\par}
    \vspace{0.8cm}
    {\LARGE Day 1: C++ 기초 문법\par}

    \vspace{2cm}

    {\large 프로그램 구조, 입출력, 변수, 연산자, 조건문, 반복문\par}

    \vfill

    {\large Summer Coding Study\par}
    \vspace{0.5cm}
    {\large \today\par}
\end{titlepage}

\tableofcontents
\newpage

% =================================================
\section{학습 목표}
% =================================================

\begin{learningbox}
이 장을 마치면 다음 내용을 설명하고 직접 코드로 작성할 수 있다.
\begin{itemize}[leftmargin=1.5em]
    \item C++ 프로그램의 기본 구조
    \item 화면 출력과 사용자 입력
    \item 변수와 기본 자료형
    \item 산술 연산자와 자동 형변환
    \item 조건문 \texttt{if}
    \item 반복문 \texttt{for}, \texttt{while}
\end{itemize}
\end{learningbox}

% =================================================
\section{C++ 프로그램의 기본 구조}
% =================================================

가장 기본적인 C++ 프로그램은 다음과 같다.

\begin{lstlisting}
#include <iostream>

using namespace std;

int main()
{
    cout << "Hello, C++!" << endl;

    return 0;
}
\end{lstlisting}

\subsection{\texttt{\#include <iostream>}}

\texttt{iostream}은 입력과 출력 기능을 제공하는 헤더이다.
\texttt{cout}, \texttt{cin}, \texttt{endl}을 사용하려면 포함해야 한다.

\subsection{\texttt{using namespace std;}}

\texttt{cout}의 실제 이름은 \texttt{std::cout}이다.
여기서 \texttt{std}는 C++ 표준 라이브러리의 이름 공간(namespace)이다.

\begin{lstlisting}
std::cout << "Hello";
\end{lstlisting}

\texttt{using namespace std;}를 작성하면 \texttt{std::}를 생략할 수 있다.

\subsection{\texttt{int main()}}

모든 C++ 프로그램은 \texttt{main()} 함수에서 시작한다.
중괄호 \texttt{\{\}}는 함수에 포함되는 코드 범위를 나타낸다.

\subsection{\texttt{return 0;}}

프로그램이 정상적으로 종료되었다는 상태 코드 0을 운영체제에 반환한다.

\begin{importantbox}
\texttt{int}, \texttt{if}, \texttt{for}, \texttt{return} 등은 C++ 언어 자체의 문법이다.
반면 \texttt{cout}, \texttt{cin}, \texttt{string}, \texttt{vector} 등은 표준 라이브러리에서 제공된다.
\end{importantbox}

% =================================================
\section{출력: \texttt{cout}}
% =================================================

출력의 기본 형태는 다음과 같다.

\begin{lstlisting}
cout << 출력할_데이터;
\end{lstlisting}

\texttt{<<}는 오른쪽 데이터를 왼쪽의 출력 스트림으로 전달한다.

\begin{lstlisting}
cout << "Age : " << 20 << endl;
\end{lstlisting}

출력 결과:

\begin{verbatim}
Age : 20
\end{verbatim}

\subsection{\texttt{endl}과 줄바꿈}

\begin{lstlisting}
cout << "Hello" << endl;
cout << "World" << endl;
\end{lstlisting}

\texttt{endl}은 줄바꿈을 수행하고 출력 버퍼도 비운다.
단순 줄바꿈에는 \texttt{\textbackslash n}도 사용할 수 있다.

\begin{lstlisting}
cout << "Hello\n";
cout << "World\n";
\end{lstlisting}

\subsection{예제: 학생 정보 출력}

\begin{lstlisting}
#include <iostream>

using namespace std;

int main()
{
    cout << "========== Student Info ==========" << endl;
    cout << "Name  : Alice" << endl;
    cout << "Major : Physics" << endl;
    cout << "Grade : 2" << endl;
    cout << "GPA   : 4.00" << endl;
    cout << "==================================" << endl;

    return 0;
}
\end{lstlisting}

% =================================================
\section{변수와 자료형}
% =================================================

변수는 \textbf{값이 저장된 메모리 공간에 붙인 이름}이다.

\begin{lstlisting}
int age = 20;
\end{lstlisting}

위 코드는 정수를 저장할 메모리 공간을 만들고, 그 공간에 \texttt{age}라는 이름을 붙인 뒤 20을 저장한다.

\subsection{기본 자료형}

\begin{center}
\begin{tabular}{lll}
\toprule
자료형 & 의미 & 예시 \\
\midrule
\texttt{int} & 정수 & \texttt{20} \\
\texttt{double} & 실수 & \texttt{3.14} \\
\texttt{char} & 문자 한 개 & \texttt{'A'} \\
\texttt{bool} & 참/거짓 & \texttt{true} \\
\texttt{string} & 문자열 & \texttt{"Hello"} \\
\bottomrule
\end{tabular}
\end{center}

문자열을 사용할 때는 \texttt{<string>} 헤더를 포함하는 것이 원칙이다.

\begin{lstlisting}
#include <iostream>
#include <string>

using namespace std;

int main()
{
    string name = "Alice";
    int age = 20;
    double gpa = 4.00;
    char grade = 'A';
    bool pass = true;

    cout << "Name  : " << name << endl;
    cout << "Age   : " << age << endl;
    cout << "GPA   : " << gpa << endl;
    cout << "Grade : " << grade << endl;
    cout << "Pass  : " << pass << endl;

    return 0;
}
\end{lstlisting}

\begin{importantbox}
\texttt{cout << 20;}은 값 20을 직접 출력한다.
\texttt{cout << age;}는 메모리에서 \texttt{age}에 저장된 값을 읽어서 출력한다.
\end{importantbox}

% =================================================
\section{입력: \texttt{cin}}
% =================================================

입력의 기본 형태는 다음과 같다.

\begin{lstlisting}
cin >> 변수;
\end{lstlisting}

\begin{lstlisting}
int age;

cout << "Enter age: ";
cin >> age;

cout << "Age : " << age << endl;
\end{lstlisting}

\subsection{여러 값 입력}

\begin{lstlisting}
int a;
int b;

cin >> a >> b;
\end{lstlisting}

다음 입력은 모두 동일하게 처리된다.

\begin{verbatim}
10 20
\end{verbatim}

또는

\begin{verbatim}
10
20
\end{verbatim}

\texttt{cin}은 공백, 탭, 줄바꿈을 입력 구분자로 사용한다.

\subsection{문자열 입력의 주의점}

\begin{lstlisting}
string name;
cin >> name;
\end{lstlisting}

사용자가 \texttt{Hong Gil Dong}을 입력하면 \texttt{name}에는 \texttt{Hong}만 저장된다.
공백을 포함한 문자열 전체는 이후에 \texttt{getline()}으로 입력받는다.

% =================================================
\section{산술 연산자}
% =================================================

\begin{center}
\begin{tabular}{cll}
\toprule
연산자 & 의미 & 예시 \\
\midrule
\texttt{+} & 덧셈 & \texttt{a + b} \\
\texttt{-} & 뺄셈 & \texttt{a - b} \\
\texttt{*} & 곱셈 & \texttt{a * b} \\
\texttt{/} & 나눗셈 & \texttt{a / b} \\
\texttt{\%} & 나머지 & \texttt{a \% b} \\
\bottomrule
\end{tabular}
\end{center}

\subsection{정수 나눗셈}

\begin{lstlisting}
int a = 7;
int b = 3;

cout << a / b << endl;
\end{lstlisting}

결과는 2이다. 두 피연산자가 모두 정수이므로 소수점 이하가 버려진다.

\subsection{실수 나눗셈과 자동 형변환}

\begin{lstlisting}
int a = 7;
double b = 3;

cout << a / b << endl;
\end{lstlisting}

이 경우 \texttt{a}가 자동으로 \texttt{double}로 변환되어 결과는 약 2.33333이 된다.

\subsection{복합 대입과 증감}

\begin{lstlisting}
x += 3;   // x = x + 3
x -= 2;   // x = x - 2
x *= 5;   // x = x * 5
x /= 2;   // x = x / 2

x++;      // x를 1 증가
x--;      // x를 1 감소
\end{lstlisting}

% =================================================
\section{조건문: \texttt{if}}
% =================================================

조건문은 조건의 참과 거짓에 따라 서로 다른 코드를 실행한다.

\begin{lstlisting}
if (조건)
{
    실행할 코드;
}
\end{lstlisting}

\subsection{비교 연산자}

\begin{center}
\begin{tabular}{cl}
\toprule
연산자 & 의미 \\
\midrule
\texttt{==} & 같다 \\
\texttt{!=} & 다르다 \\
\texttt{>} & 크다 \\
\texttt{<} & 작다 \\
\texttt{>=} & 크거나 같다 \\
\texttt{<=} & 작거나 같다 \\
\bottomrule
\end{tabular}
\end{center}

\begin{importantbox}
\texttt{=}는 값을 대입하는 연산자이고, \texttt{==}는 두 값이 같은지 비교하는 연산자이다.
\end{importantbox}

\subsection{\texttt{if-else}}

\begin{lstlisting}
if (score >= 60)
{
    cout << "Pass" << endl;
}
else
{
    cout << "Fail" << endl;
}
\end{lstlisting}

\subsection{\texttt{else if}}

\begin{lstlisting}
if (score >= 90)
{
    cout << "A";
}
else if (score >= 80)
{
    cout << "B";
}
else if (score >= 70)
{
    cout << "C";
}
else
{
    cout << "F";
}
\end{lstlisting}

위에서부터 조건을 검사하며, 처음으로 참이 된 블록 하나만 실행한다.

\subsection{논리 연산자}

\begin{center}
\begin{tabular}{cl}
\toprule
연산자 & 의미 \\
\midrule
\texttt{\&\&} & AND, 두 조건이 모두 참 \\
\texttt{||} & OR, 둘 중 하나 이상 참 \\
\texttt{!} & NOT, 참과 거짓을 반전 \\
\bottomrule
\end{tabular}
\end{center}

% =================================================
\section{반복문: \texttt{for}}
% =================================================

반복 횟수를 알고 있을 때 일반적으로 \texttt{for}문을 사용한다.

\begin{lstlisting}
for (초기식; 조건식; 증감식)
{
    반복할 코드;
}
\end{lstlisting}

\begin{lstlisting}
for (int i = 1; i <= 5; i++)
{
    cout << i << endl;
}
\end{lstlisting}

실행 순서는 다음과 같다.

\begin{enumerate}
    \item 초기식 실행
    \item 조건식 검사
    \item 코드 블록 실행
    \item 증감식 실행
    \item 조건식으로 돌아감
\end{enumerate}

\subsection{누적}

\begin{lstlisting}
int sum = 0;

for (int i = 1; i <= 10; i++)
{
    sum += i;
}

cout << sum << endl;
\end{lstlisting}

\texttt{sum}에 이전 결과를 계속 더하는 방식을 누적(accumulation)이라고 한다.

\begin{checkpointbox}
다음 반복문에서 마지막으로 출력되는 값은 4이지만, 조건 검사를 실패할 때의 \texttt{i} 값은 5이다.
\begin{lstlisting}
for (int i = 0; i < 5; i++)
{
    cout << i << endl;
}
\end{lstlisting}
\end{checkpointbox}

% =================================================
\section{반복문: \texttt{while}}
% =================================================

\texttt{while}문은 조건이 참인 동안 반복한다.

\begin{lstlisting}
while (조건)
{
    반복할 코드;
}
\end{lstlisting}

\begin{lstlisting}
int i = 1;

while (i <= 5)
{
    cout << i << endl;
    i++;
}
\end{lstlisting}

\subsection{\texttt{for}와 \texttt{while}의 선택}

\begin{center}
\begin{tabular}{ll}
\toprule
상황 & 추천 반복문 \\
\midrule
반복 횟수를 알고 있음 & \texttt{for} \\
반복 횟수를 미리 알 수 없음 & \texttt{while} \\
\bottomrule
\end{tabular}
\end{center}

예를 들어 1부터 100까지 출력할 때는 \texttt{for}가 자연스럽고,
비밀번호를 맞힐 때까지 입력받는 경우에는 \texttt{while}이 자연스럽다.

\subsection{무한 반복}

\begin{lstlisting}
int i = 1;

while (i <= 5)
{
    cout << i << endl;
}
\end{lstlisting}

위 코드에서는 \texttt{i}가 변하지 않으므로 조건이 계속 참이다.
따라서 반복문이 끝나지 않는다.

\subsection{비밀번호 입력 예제}

\begin{lstlisting}
#include <iostream>

using namespace std;

int main()
{
    int true_password = 1234;
    int input_password = 0;

    while (input_password != true_password)
    {
        cout << "Enter password: ";
        cin >> input_password;

        if (input_password == true_password)
        {
            cout << "Access Granted" << endl;
        }
        else
        {
            cout << "Wrong Password" << endl;
            cout << endl;
        }
    }

    return 0;
}
\end{lstlisting}

\texttt{while}은 조건을 먼저 검사하므로, 위 구조에서는
\texttt{input\_password}에 초기값을 주어야 한다.
초기화하지 않은 지역 변수에는 예측할 수 없는 값이 들어 있을 수 있다.

% =================================================
\section{실습 목록}
% =================================================

\begin{practicebox}
Day 1 실습 파일은 다음과 같이 관리한다.
\begin{itemize}[leftmargin=1.5em]
    \item \texttt{practice/p01\_output.cpp}
    \item \texttt{practice/p02\_variables.cpp}
    \item \texttt{practice/p03\_rectangle\_input.cpp}
    \item \texttt{practice/p04\_rectangle\_area.cpp}
    \item \texttt{practice/p05\_even\_odd.cpp}
    \item \texttt{practice/p06\_sum.cpp}
    \item \texttt{practice/p07\_password.cpp}
\end{itemize}
\end{practicebox}

% =================================================
\section{핵심 요약}
% =================================================

\begin{enumerate}
    \item C++ 프로그램은 \texttt{main()} 함수에서 시작한다.
    \item \texttt{cout}은 출력, \texttt{cin}은 입력에 사용한다.
    \item 변수는 메모리 공간에 붙인 이름이다.
    \item \texttt{int / int}의 결과는 정수이며 소수점 이하는 버려진다.
    \item 피연산자 중 하나가 \texttt{double}이면 자동 형변환이 발생할 수 있다.
    \item \texttt{if-else if-else}에서는 첫 번째로 참인 블록만 실행된다.
    \item 반복 횟수를 알면 \texttt{for}, 모르면 \texttt{while}이 일반적으로 자연스럽다.
    \item 반복문에서는 조건을 변화시키는 코드가 없으면 무한 반복이 발생할 수 있다.
\end{enumerate}

\end{document}
